%% kapa: exam=1 round=22 year=2011 period=2 form=A subject=accounting subject_ko=회계학 seq=39 number=39
%% source: Q-Net 감정평가사 기출문제 (cst00309) artlSeq=5036790 file=extracted/1008/5036790/2011년 제22회 감정평가사 1차시험 기출문제/2교시.hwp
%% conversion: kordoc 4.24.51 -> tools/gen1.py (auto); review=proofread-sonnet; flags=none
\begin{question}{39}
다음은 상품매매업을 영위하는 (주)서울의 20×1년 예산자료 중, 1분기의 매입 및 매출 추정액에 관한 자료이다.

{\small\centering
\begin{tabular}{|c|r|r|}
\hline
\uline{월} & \multicolumn{1}{c|}{\uline{매입추정액}} & \multicolumn{1}{c|}{\uline{매출추정액}} \\
1 & ￦400,000 & ￦620,000 \\
2 & 900,000 & 748,800 \\
3 & 500,000 & 780,000 \\
\hline
\end{tabular}\par}

(주)서울의 20×1년 3월초 예상되는 현금보유액은 ￦50,000이다. 매출채권은 당월에 70\%, 다음 달에 25\%를 회수하고, 나머지는 회수하지 못할 것으로 예상된다. 매입채무는 매입한 달에 60\%를 지급하는데 그 지급액의 2\%는 현금할인 혜택을 받으며, 나머지는 다음 달에 모두 지급한다. 3월 중에 일반경비로 ￦200,000을 현금지출 할 예정이다. (주)서울이 매월 말 현금 ￦50,000을 보유하고 있어야 한다면, 3월에 추가로 조달해야 할 현금은 얼마로 추정되는가? \par
\begin{choices}
  \item ￦70,800
  \item ￦80,800
  \item ￦90,800
  \item ￦110,800
  \item ￦120,800
\end{choices}
\end{question}
\answer{5}
